\documentclass[12pt,a4paper]{article}
\usepackage[utf8]{inputenc}
\usepackage[margin=0.8in]{geometry}
\usepackage{graphicx}
\usepackage{array}
\usepackage{longtable}
\usepackage{booktabs}
\usepackage{xcolor}
\usepackage{hyperref}
\usepackage{listings}
\usepackage{fancyhdr}

% Header and Footer
\pagestyle{fancy}
\fancyhf{}
\rhead{Lab Evaluation-1}
\lhead{React JS - StellarStep}
\rfoot{Page \thepage}

% Code listing style
\lstset{
    basicstyle=\ttfamily\footnotesize,
    breaklines=true,
    frame=single,
    backgroundcolor=\color{gray!10}
}

\begin{document}

% Title Page
\begin{titlepage}
    \centering
    \vspace*{1.5cm}
    
    {\Huge\bfseries Lab Evaluation-1\par}
    \vspace{0.8cm}
    {\Large Date: 5/2/2026\par}
    \vspace{1.5cm}
    
    {\LARGE\textbf{React JS Application Development}\par}
    \vspace{0.4cm}
    {\Large StellarStep - The Sensory Space Odyssey™\par}
    \vspace{2.5cm}
    
    {\Large
    \textbf{Name:} Rohith Kumar D\par}
    \vspace{0.3cm}
    {\Large
    \textbf{Roll No:} cb.sc.u4cse23018\par}
    \vspace{0.3cm}
    {\Large
    \textbf{Section:} CSE-A\par}
    \vspace{2.5cm}
    
    {\large\textbf{Total Marks: 50}\par}
    \vspace{0.2cm}
    {\normalsize 60\% weightage for earlier submission and 40\% for today's evaluation\par}
    
    \vfill
    
\end{titlepage}

% Main Content
\section{About the Use Case}

StellarStep - The Sensory Space Odyssey™ is a comprehensive neuro-inclusive web application specifically designed for children with Autism Spectrum Disorder (ASD). The application provides an immersive space-themed learning environment that combines sensory-friendly design principles with interactive educational content. The primary objective is to create an accessible, engaging, and therapeutic digital space where children with ASD can learn, explore, and develop essential skills through gamification and sensory experiences.

The application features voice navigation with full command support using Web Speech API, sensory regulation through calming nebula environments with interactive particle effects, and multiple educational missions including Planet Matcher, Alien Emotions, and Space School. It includes progress tracking via visual timeline and stellar gallery, customization options through the dress-up station, and a Focus Trainer for attention development. The space theme provides clear metaphors for abstract concepts (emotions as alien faces, time management as space missions), making complex ideas more accessible.

Built with modern technologies including React 18.2.0, Framer Motion for animations, TailwindCSS for styling, and Vite as the build tool, the application targets children aged 8-14 with ASD, their parents, caregivers, special education teachers, therapists, and educational institutions focused on inclusive learning. The design emphasizes accessibility with large touch targets, clear visual hierarchy, and consistent design patterns throughout.

\section{Extended Lab Exercise - Space Lab Page}

As an extension to Lab 1, a comprehensive new page called "Space Lab" has been implemented, demonstrating all React concepts including advanced hooks, class components, forms with validation, and state management patterns. The page features an Astronaut Registration Form built as a Class Component with full lifecycle methods, a Mission Planning System using the useReducer hook for advanced state management, a Statistics Dashboard with performance-optimized calculations using useMemo, and interactive components combining stateful and stateless components with PropTypes validation. The Space Lab page is integrated into the main application navigation system, accessible via the header menu with a 🔬 icon and voice commands ("Space Lab" or "Lab").

\newpage

\section{React JS Concepts Implementation}

This section provides comprehensive documentation of all React JS concepts implemented in the StellarStep application, with specific focus on the Space Lab extension page.

\vspace{1cm}

\begin{longtable}{|p{3.5cm}|p{6cm}|p{5cm}|}
\hline
\textbf{Concept} & \textbf{Code Proof / Location} & \textbf{Screenshot} \\
\hline
\endfirsthead

\hline
\textbf{Concept} & \textbf{Code Proof / Location} & \textbf{Screenshot} \\
\hline
\endhead

\hline
\endfoot

\hline
\endlastfoot

\textbf{Function Component} & 
\textbf{File:} \texttt{SpaceLab.jsx} (Line 404)
\newline\newline
\textbf{Implementation:}
\begin{lstlisting}[language=JavaScript]
const SpaceLab = () => {
  const { speak } = useVoice();
  const [astronauts, 
    setAstronauts] = useState([]);
  
  return (
    <div className="min-h-screen">
      {/* Page content */}
    </div>
  );
};
\end{lstlisting}
\newline
\textbf{Also Implemented:}
\begin{itemize}
    \item MissionCard (Line 346)
    \item All page components
    \item Navbar, LoadingScreen
\end{itemize}
& 
\includegraphics[width=4.5cm,height=6cm,keepaspectratio]{figure-1.png}
\newline
\textit{Figure 1: Function Component Implementation}
\\
\hline

\textbf{Class Component} & 
\textbf{File:} \texttt{SpaceLab.jsx} (Lines 24-332)
\newline\newline
\textbf{Class Definition:}
\begin{lstlisting}[language=JavaScript]
class AstronautRegistrationForm 
  extends React.Component {
  
  constructor(props) {
    super(props);
    this.state = {
      name: '',
      age: '',
      rank: 'cadet',
      skills: [],
      bio: '',
      errors: {},
      submitted: false
    };
  }
  
  componentDidMount() {
    console.log('Form mounted');
  }
  
  componentDidUpdate(prevProps, 
    prevState) {
    if (prevState.submitted !== 
      this.state.submitted) {
      console.log('Submitted');
    }
  }
  
  componentWillUnmount() {
    console.log('Unmounting');
  }
  
  render() {
    return (/* JSX */);
  }
}
\end{lstlisting}
& 
\includegraphics[width=4.5cm,height=6cm,keepaspectratio]{figure-2.png}
\newline
\textit{Figure 2: Class Component with Lifecycle Methods}
\\
\hline

\textbf{Event Handling} & 
\textbf{File:} \texttt{SpaceLab.jsx}
\newline\newline
\textbf{Event Types Implemented:}
\begin{lstlisting}[language=JavaScript]
// onChange event
<input
  type="text"
  name="name"
  value={name}
  onChange={this.handleInputChange}
/>

// onSubmit event
<form onSubmit={this.handleSubmit}>
  {/* form fields */}
</form>

// onClick event
<button
  onClick={handleDeleteAstronaut}
>
  Delete
</button>

// Checkbox onChange
<input
  type="checkbox"
  value={skill}
  checked={skills.includes(skill)}
  onChange={this.handleSkillChange}
/>
\end{lstlisting}
\newline
\textbf{Locations:}
\begin{itemize}
    \item onChange: Lines 119, 221, 229
    \item onSubmit: Line 153
    \item onClick: Lines 271, 456
\end{itemize}
& 
\includegraphics[width=4.5cm,height=6cm,keepaspectratio]{figure-3.png}
\newline
\textit{Figure 3: Event Handlers in Action}
\\
\hline

\textbf{State Management} & 
\textbf{Multiple Patterns Implemented:}
\newline\newline
\textbf{1. useState Hook:}
\begin{lstlisting}[language=JavaScript]
const [astronauts, 
  setAstronauts] = useState([]);
const [newMission, setNewMission] 
  = useState({ 
    title: '', 
    description: '' 
  });
\end{lstlisting}
\newline
\textbf{2. useReducer Hook:}
\begin{lstlisting}[language=JavaScript]
const missionReducer = (state, 
  action) => {
  switch (action.type) {
    case 'ADD_MISSION':
      return {
        ...state,
        missions: [
          ...state.missions, 
          action.payload
        ]
      };
    case 'DELETE_MISSION':
      return { /* ... */ };
    default:
      return state;
  }
};

const [missionState, dispatch] = 
  useReducer(missionReducer, {
    missions: [],
    totalMissions: 0
  });
\end{lstlisting}
\newline
\textbf{3. Class Component State:}
\begin{lstlisting}[language=JavaScript]
this.state = {
  name: '',
  errors: {}
};

this.setState({ 
  name: value 
});
\end{lstlisting}
& 
\includegraphics[width=4.5cm,height=6cm,keepaspectratio]{figure-4.png}
\newline
\textit{Figure 4: State Management with useState and useReducer}
\\
\hline

\textbf{Stateless Component} & 
\textbf{File:} \texttt{SpaceLab.jsx} (Lines 346-377)
\newline\newline
\textbf{Pure Presentational Component:}
\begin{lstlisting}[language=JavaScript]
const MissionCard = ({ 
  mission, 
  onToggle, 
  onDelete 
}) => {
  return (
    <motion.div>
      <div className="flex items-center">
        <input
          type="checkbox"
          checked={mission.completed}
          onChange={() => 
            onToggle(mission.id)}
        />
        <h3>{mission.title}</h3>
        <button 
          onClick={() => 
            onDelete(mission.id)}
        >
          Delete
        </button>
      </div>
    </motion.div>
  );
};

MissionCard.propTypes = {
  mission: PropTypes.shape({
    id: PropTypes.number.isRequired,
    title: PropTypes.string.isRequired,
    completed: PropTypes.bool.isRequired
  }).isRequired,
  onToggle: PropTypes.func.isRequired,
  onDelete: PropTypes.func.isRequired
};
\end{lstlisting}
\newline
\textbf{Characteristics:}
\begin{itemize}
    \item No internal state
    \item Receives all data via props
    \item PropTypes for validation
    \item Pure presentation logic
\end{itemize}
& 
\includegraphics[width=4.5cm,height=6cm,keepaspectratio]{figure-5.png}
\newline
\textit{Figure 5: Stateless MissionCard Component}
\\
\hline

\textbf{Forms with Validation} & 
\textbf{File:} \texttt{SpaceLab.jsx} (Lines 153-291)
\newline\newline
\textbf{Complete Form Implementation:}
\begin{lstlisting}[language=JavaScript]
<form onSubmit={this.handleSubmit}>
  {/* Text Input */}
  <input
    type="text"
    name="name"
    value={name}
    onChange={this.handleInputChange}
    placeholder="Astronaut name"
  />
  {errors.name && <p>{errors.name}</p>}
  
  {/* Number Input */}
  <input
    type="number"
    name="age"
    value={age}
    onChange={this.handleInputChange}
    min="8"
    max="100"
  />
  
  {/* Select Dropdown */}
  <select
    name="rank"
    value={rank}
    onChange={this.handleInputChange}
  >
    <option value="cadet">Cadet</option>
    <option value="pilot">Pilot</option>
    <option value="commander">
      Commander
    </option>
  </select>
  
  {/* Checkboxes */}
  {['Piloting', 'Engineering'].map(
    skill => (
    <input
      type="checkbox"
      value={skill}
      checked={skills.includes(skill)}
      onChange={this.handleSkillChange}
    />
  ))}
  
  {/* Textarea */}
  <textarea
    name="bio"
    value={bio}
    onChange={this.handleInputChange}
    rows="4"
  />
  
  <button type="submit">
    Register
  </button>
  <button type="button" 
    onClick={this.handleReset}>
    Reset
  </button>
</form>
\end{lstlisting}
\newline
\textbf{Validation Logic:}
\begin{lstlisting}[language=JavaScript]
validateForm = () => {
  const errors = {};
  if (!name.trim()) {
    errors.name = 'Name required';
  }
  if (age < 8 || age > 100) {
    errors.age = 'Age 8-100';
  }
  if (skills.length === 0) {
    errors.skills = 'Select skill';
  }
  return errors;
};
\end{lstlisting}
& 
\includegraphics[width=4.5cm,height=6cm,keepaspectratio]{figure-6.png}
\newline
\textit{Figure 6: Complete Form with Validation and Multiple Input Types}
\\
\hline

\end{longtable}

\newpage

\section{Additional React Concepts Implemented}

\subsection{Advanced Hooks}

\subsubsection{useReducer Hook}
Implemented in Mission Planning system for complex state management with actions like ADD\_MISSION, DELETE\_MISSION, TOGGLE\_COMPLETE, and RESET.

\subsubsection{useMemo Hook}
Used for performance optimization in statistics calculations:
\begin{lstlisting}[language=JavaScript]
const stats = useMemo(() => {
  const totalAstronauts = astronauts.length;
  const completedMissions = missionState.missions
    .filter(m => m.completed).length;
  const avgAge = astronauts.reduce((sum, a) => 
    sum + parseInt(a.age), 0) / astronauts.length;
  return { totalAstronauts, completedMissions, avgAge };
}, [astronauts, missionState.missions]);
\end{lstlisting}

\subsubsection{useCallback Hook}
Memoized event handlers to prevent unnecessary re-renders:
\begin{lstlisting}[language=JavaScript]
const handleAstronautSubmit = useCallback((astronaut) => {
  setAstronauts(prev => [...prev, astronaut]);
  speak(`${astronaut.name} registered`);
}, [speak]);
\end{lstlisting}

\subsubsection{useEffect Hook}
Lifecycle management and side effects:
\begin{lstlisting}[language=JavaScript]
useEffect(() => {
  speak('Welcome to Space Lab!');
}, []);
\end{lstlisting}

\subsubsection{useContext Hook}
Custom hook for accessing VoiceContext:
\begin{lstlisting}[language=JavaScript]
export const useVoice = () => {
  const context = useContext(VoiceContext);
  if (!context) {
    throw new Error('useVoice must be used within VoiceProvider');
  }
  return context;
};
\end{lstlisting}

\subsection{PropTypes}
Type checking implemented on multiple components:
\begin{lstlisting}[language=JavaScript]
AstronautRegistrationForm.propTypes = {
  onSubmit: PropTypes.func.isRequired,
  onMount: PropTypes.func,
  theme: PropTypes.shape({
    primary: PropTypes.string,
    secondary: PropTypes.string
  })
};
\end{lstlisting}

\subsection{Higher-Order Component (HOC)}
Theme wrapper HOC implementation:
\begin{lstlisting}[language=JavaScript]
const withSpaceTheme = (Component) => {
  return function ThemedComponent(props) {
    const theme = {
      primary: 'cosmic-purple',
      secondary: 'cosmic-blue'
    };
    return <Component {...props} theme={theme} />;
  };
};
\end{lstlisting}

\subsection{React.Fragment}
Used in Statistics Dashboard to avoid unnecessary DOM nodes:
\begin{lstlisting}[language=JavaScript]
<Fragment>
  <div>Total Astronauts: {stats.totalAstronauts}</div>
  <div>Completed: {stats.completedMissions}</div>
  <div>Pending: {stats.pendingMissions}</div>
  <div>Avg Age: {stats.avgAge}</div>
</Fragment>
\end{lstlisting}

\subsection{Context API}
Global state management for voice navigation:
\begin{lstlisting}[language=JavaScript]
const VoiceContext = createContext();

export const VoiceProvider = ({ children }) => {
  const [isVoiceEnabled, setIsVoiceEnabled] = useState(true);
  const value = { isVoiceEnabled, setIsVoiceEnabled, speak };
  
  return (
    <VoiceContext.Provider value={value}>
      {children}
    </VoiceContext.Provider>
  );
};
\end{lstlisting}

\subsection{Custom Navigation System}
State-based routing implementation without React Router:
\begin{lstlisting}[language=JavaScript]
const [currentPage, setCurrentPage] = useState('landing');

const navigateToPage = (pageName) => {
  setCurrentPage(pageName);
};

const pages = {
  'landing': <LandingPage />,
  'mission-hub': <MissionHub />,
  'space-lab': <SpaceLab />
};

return pages[currentPage];
\end{lstlisting}

\newpage

\section{Evaluation Criteria}

\subsection{UI Aspects (10 Marks)}
\begin{itemize}
    \item \textbf{Design Consistency:} Space theme maintained throughout with consistent color scheme
    \item \textbf{Responsiveness:} Fully responsive design using Tailwind CSS grid and flexbox
    \item \textbf{Animations:} Smooth transitions using Framer Motion for enhanced UX
    \item \textbf{Accessibility:} Large touch targets (minimum 56px height), clear visual hierarchy
    \item \textbf{User Feedback:} Visual indicators for form validation, success messages, loading states
    \item \textbf{Typography:} Custom 'font-friendly' class for readability
    \item \textbf{Color Contrast:} High contrast ratios for text visibility
    \item \textbf{Interactive Elements:} Hover effects, click animations, focus states
\end{itemize}

\subsection{React JS Concepts (30 Marks)}
\begin{itemize}
    \item \textbf{Function Components:} SpaceLab, MissionCard, all page components
    \item \textbf{Class Component:} AstronautRegistrationForm with full lifecycle
    \item \textbf{Hooks:} useState, useEffect, useContext, useReducer, useMemo, useCallback, useRef
    \item \textbf{State Management:} Local state, reducer pattern, Context API
    \item \textbf{Forms:} Complex form with text, number, select, checkbox, textarea inputs
    \item \textbf{Validation:} Client-side validation with error messages
    \item \textbf{Events:} onChange, onSubmit, onClick, onMouseEnter handlers
    \item \textbf{PropTypes:} Type checking on stateless components
    \item \textbf{HOC:} withSpaceTheme higher-order component
    \item \textbf{Fragments:} Used for clean DOM structure
    \item \textbf{Conditional Rendering:} Multiple patterns throughout
    \item \textbf{Lists \& Keys:} Proper key usage in mapped arrays
\end{itemize}

\subsection{Extension Level (10 Marks)}
\begin{itemize}
    \item \textbf{Advanced Features:} Voice navigation using Web Speech API
    \item \textbf{Performance Optimization:} useMemo for expensive calculations
    \item \textbf{Code Organization:} Separate folders for components, pages, contexts
    \item \textbf{Reusability:} Modular component structure with props
    \item \textbf{Error Handling:} Form validation and error boundaries
    \item \textbf{State Persistence:} localStorage for user preferences
    \item \textbf{Complex State Logic:} useReducer for mission management
    \item \textbf{Type Safety:} PropTypes implementation
    \item \textbf{Accessibility Features:} Voice commands, keyboard navigation
    \item \textbf{Animation Integration:} Framer Motion for smooth transitions
\end{itemize}

\newpage

\section{Project Structure}

\subsection{Directory Organization}
\begin{verbatim}
stellar_frontend_application/
├── src/
│   ├── components/
│   │   ├── Navbar.jsx
│   │   ├── MusicPlayer.jsx
│   │   └── LoadingScreen.jsx
│   ├── contexts/
│   │   └── VoiceContext.jsx
│   ├── pages/
│   │   ├── LandingPage.jsx
│   │   ├── MissionHub.jsx
│   │   ├── PlanetMatcher.jsx
│   │   ├── AlienEmotions.jsx
│   │   ├── SensoryNebula.jsx
│   │   ├── StellarGallery.jsx
│   │   ├── SpaceSchool.jsx
│   │   ├── FocusTrainer.jsx
│   │   ├── DressUpStation.jsx
│   │   ├── VisualTimeline.jsx
│   │   └── SpaceLab.jsx (NEW)
│   ├── assets/
│   ├── App.jsx
│   ├── main.jsx
│   └── index.css
├── public/
│   └── assets/
├── package.json
├── vite.config.js
├── tailwind.config.js
└── README.md
\end{verbatim}

\subsection{Component Hierarchy}
\begin{itemize}
    \item \textbf{App.jsx} - Root component with navigation logic
    \item \textbf{VoiceProvider} - Context provider wrapping the entire app
    \item \textbf{Navbar} - Global navigation component
    \item \textbf{Pages} - Individual route components
    \item \textbf{SpaceLab} - Extended page with:
    \begin{itemize}
        \item AstronautRegistrationForm (Class Component)
        \item MissionCard (Stateless Component)
        \item Statistics Dashboard (useMemo)
        \item Mission Planning System (useReducer)
    \end{itemize}
\end{itemize}

\newpage

\section{Technical Implementation Details}

\subsection{Dependencies}
\begin{verbatim}
{
  "dependencies": {
    "react": "^18.2.0",
    "react-dom": "^18.2.0",
    "framer-motion": "^10.16.0",
    "lucide-react": "^0.294.0",
    "prop-types": "^15.8.1"
  },
  "devDependencies": {
    "@vitejs/plugin-react": "^4.2.0",
    "autoprefixer": "^10.4.16",
    "postcss": "^8.4.32",
    "tailwindcss": "^3.3.6",
    "vite": "^5.0.7"
  }
}
\end{verbatim}

\subsection{Build and Development}
\begin{itemize}
    \item \textbf{Build Tool:} Vite 5.0.7 for fast development and optimized builds
    \item \textbf{Development Server:} Hot Module Replacement (HMR) enabled
    \item \textbf{CSS Framework:} Tailwind CSS with custom configuration
    \item \textbf{Animation Library:} Framer Motion for smooth transitions
    \item \textbf{Icons:} Lucide React icon library
\end{itemize}

\subsection{Performance Optimizations}
\begin{itemize}
    \item useMemo for expensive statistics calculations
    \item useCallback for memoized event handlers
    \item Code splitting with dynamic imports
    \item Lazy loading of components
    \item Optimized re-renders with proper dependency arrays
\end{itemize}

\newpage

\section{Testing and Validation}

\subsection{Browser Compatibility}
\begin{itemize}
    \item Chrome 90+
    \item Firefox 88+
    \item Safari 14+
    \item Edge 90+
\end{itemize}

\subsection{Feature Testing}
\begin{itemize}
    \item \textbf{Form Validation:} All input fields validated with appropriate error messages
    \item \textbf{State Management:} useState and useReducer tested with multiple actions
    \item \textbf{Voice Navigation:} Commands tested across all pages
    \item \textbf{Responsive Design:} Tested on mobile, tablet, and desktop viewports
    \item \textbf{Accessibility:} Keyboard navigation and screen reader compatibility
\end{itemize}

\subsection{Code Quality}
\begin{itemize}
    \item PropTypes validation on all components with props
    \item Consistent code formatting and naming conventions
    \item Proper error handling and boundary conditions
    \item Clean component separation and single responsibility principle
    \item Comprehensive comments and documentation
\end{itemize}

\newpage

\section{Conclusion}

\subsection{Learning Outcomes}
Through the development of StellarStep and the Space Lab extension, the following React JS concepts have been successfully implemented and demonstrated:

\begin{enumerate}
    \item \textbf{Component Architecture:} Both functional and class-based components with proper lifecycle management
    \item \textbf{Modern Hooks:} All major React hooks including useState, useEffect, useContext, useReducer, useMemo, and useCallback
    \item \textbf{Form Handling:} Complex forms with validation, multiple input types, and user feedback
    \item \textbf{State Management:} Multiple patterns including local state, reducer pattern, and Context API
    \item \textbf{Performance:} Optimization techniques using memoization and proper dependency management
    \item \textbf{Type Safety:} PropTypes implementation for runtime type checking
    \item \textbf{Advanced Patterns:} Higher-Order Components and custom hooks
\end{enumerate}

\subsection{Future Enhancements}
Potential improvements and extensions include:
\begin{itemize}
    \item React Router implementation for URL-based navigation
    \item TypeScript migration for compile-time type safety
    \item Unit testing with Jest and React Testing Library
    \item Backend integration with REST API or GraphQL
    \item Progressive Web App (PWA) features
    \item Enhanced accessibility features (ARIA labels, screen reader optimization)
\end{itemize}

\subsection{Acknowledgments}
This project demonstrates comprehensive understanding of React JS fundamentals and advanced concepts as outlined in the TutorialsPoint React.js tutorial. All required concepts have been implemented with proper code organization and best practices.

\vspace{2cm}

\begin{center}
\textbf{--- End of Document ---}
\end{center}

\end{document}
